% TOP_PROBLEM: {"id":30007621,"problem_number":"LOCAL-30007621","title":"Global financial safety net","category_id":21,"category":{"id":21,"name":"economics","display_name":"Economics","description":"Open economic theory statements, published research agendas, and proposed scoped specifications; question provenance and historical priority are qualified per record.","slug":"economics","order_index":21,"created_at":"2026-10-08T00:00:00Z"},"status":"research_question","external_url":null,"legacy_ids":[],"published":false,"statement":"In the explicitly specified setting: What combination of reserves, swap lines, regional facilities and IMF support best insures countries against systemic crises? The mathematical statement above fixes the target, assumptions and scope of this catalogue formulation.","economics":{"original_id":110,"field":"International finance and exchange rates","kind":"design","formalization_basis":"adapted_research_question","source_status":"research_agenda","priority_status":"earliest_found","priority_note":"This is the earliest explicit statement verified in this audit; first priority for the full catalogue formulation is not established.","earliest_verified_reference":{"authors":["Bank of England"],"title":"Bank of England Agenda for Research, 2025–2028","year":2026,"url":"https://www.bankofengland.co.uk/research/bank-of-england-agenda-for-research","doi":null,"locator":"Interconnected World / International finance, paragraph 1","evidence_summary":"Requests safety-net reform and coordinated crisis tools."},"current_reference":{"authors":["Bank of England"],"title":"Bank of England Agenda for Research, 2025–2028","year":2026,"url":"https://www.bankofengland.co.uk/research/bank-of-england-agenda-for-research","doi":null,"locator":"Interconnected World / International finance, paragraph 1","evidence_summary":"Requests safety-net reform and coordinated crisis tools."},"formalization_note":"Adapts an explicitly verified research-agenda passage. Mathematical objects, interventions, objectives and scope are proposed here; existing model-specific solutions are not relabeled unsolved."},"statusQualification":"Published question or research agenda verified at the cited locator. The mathematical specification is adapted for this catalogue and includes additional assumptions; source publication does not establish first-ever priority or certify this exact model remains unsolved. Adapts an explicitly verified research-agenda passage. Mathematical objects, interventions, objectives and scope are proposed here; existing model-specific solutions are not relabeled unsolved.","statusReviewedAt":"2026-10-08"}
% FIRST_POSTED: 2026-10-08
% ENTRY_KIND: statement_only
% !TeX program = lualatex
\documentclass[11pt]{article}
\usepackage{fontspec}
\usepackage[a4paper,margin=25mm]{geometry}
\usepackage{amsmath,amssymb,mathtools}
\usepackage{xurl}
\usepackage[hidelinks]{hyperref}
\newcommand{\sref}[2]{\hyperlink{source:#1:#2}{[S#2]}}
\setlength{\emergencystretch}{3em}
\begin{document}
% problemId:30007621
\section*{Global financial safety net}\label{30007621}
\subsection{Definitions and mathematical statement}
Consider countries choosing reserve buffers and accessing central-bank swaps, regional credit facilities and multilateral lending. A safety-net architecture \(a\) specifies eligibility, committed capacities, prices, collateral, conditionality and activation rules. Countries choose prior borrowing and self-insurance anticipating access; lenders face a specified resource constraint. Let \(\omega\) be a shock state with a specified probability law. Let \(L_c(a,\omega)\) be real crisis losses in country \(c\), \(C_c(a)\) reserve and facility costs, and \(F(a,\omega)\) shared fiscal resource losses. Given nonnegative welfare weights \(w_c\), seek a feasible architecture minimizing
\[\sum_cw_c\{E[L_c(a,\omega)]+C_c(a)\}+E[F(a,\omega)].\]
The task is to characterize coverage, coordination and instrument substitution that improve insurance without inducing excessive prior risk. Define stigma as a measurable private access cost or signaling consequence, and incorporate it in country behavior. Account for correlated global demands for liquidity and currency-specific funding needs. Identify causal effects where past support was selected in response to severity. Report sensitivity to unequal eligibility and alternative welfare weights. Success in protecting eligible countries is insufficient to establish optimal global coverage or resource allocation.

\par\textbf{Formulation and scope.} Adapts an explicitly verified research-agenda passage. Mathematical objects, interventions, objectives and scope are proposed here; existing model-specific solutions are not relabeled unsolved.

\par\textbf{Open-question provenance.} An explicit question or research agenda was verified at the source locator. This does not by itself certify that every later version remains unresolved \sref{30007621}{1}.

\par\textbf{Historical priority.} This is the earliest explicit statement verified in this audit; first priority for the full catalogue formulation is not established.

\subsection{Short English statement}
In the explicitly specified setting: What combination of reserves, swap lines, regional facilities and IMF support best insures countries against systemic crises? The mathematical statement above fixes the target, assumptions and scope of this catalogue formulation.

\subsection{Sources}
\begin{itemize}\raggedright
\item[\textnormal{[S1]}]\hypertarget{source:30007621:1}{} Bank of England. 2026. Bank of England Agenda for Research, 2025–2028. Interconnected World / International finance, paragraph 1.
\url{https://www.bankofengland.co.uk/research/bank-of-england-agenda-for-research}.
{\small\textit{Earliest verified question/agenda reference; historical priority qualified below. Requests safety-net reform and coordinated crisis tools.}}
\end{itemize}
\end{document}
